\section{Track F: QMC and Applications Part II}\newpage

\begin{talk}
  {Approximation of multivariate periodic functions}% [1] talk title
  {Dirk Nuyens}% [2] speaker name
  {KU Leuven, Belgium}% [3] affiliations
  {dirk.nuyens@kuleuven.be}% [4] email
  {Laurence Wilkes}% [5] coauthors
  {QMC and Applications}% [6] special session. Leave this field empty for contributed talks. 
    
   
  We study approximation of multivariate periodic functions using $n$
  function values. We use generated sets as the sample points, first
  introduced in~[1]. We prove existence and convergence of the almost
  optimal $L_2$ error and obtain similar bounds as for the least squares
  algorithm from~[2] which uses unstructured points.

\medskip

\begin{enumerate}
  \item[{[1]}]
    K{\"a}mmerer.
    Reconstructing multivariate trigonometric polynomials by sampling along generated sets.
    In Monte Carlo and Quasi-Monte Carlo Methods 2012 (Dick, Kuo, Peters, Sloan), pages 439--454, 2013.
  \item[{[2]}]
    Krieg, Ullrich.
    Function values are enough for $L_2$ approximation.
    \emph{Foundations of Computational Mathematics}, 21(4):1141--1151, 2021.
\end{enumerate}

\end{talk}

\begin{talk}
  {Randomized QMC with one categorical variable}% [1] talk title
  {Art B. Owen}% [2] speaker name
  {Stanford University}% [3] affiliations
  {owen@stanford.edu}% [4] email
  {Valerie Ho, Zexin Pan}% [5] coauthors
  
    
   
  Randomized quasi-Monte Carlo (RQMC) methods benefit from smoothness
  in the integrand.  In some applications, one of the input variables takes
  only a modest finite number $L\geqslant2$
of values and can be considered a categorical
  variable. Such a variable introduces a discontinuity in the integrand.  It is
  an RQMC-friendly discontinuity because the discontinuity is axis parallel
  when each level of the categorical variable corresponds to a single
  sub-interval of $[0,1]$.  A naturally occurring use case has
  the categorical variable correspond to a component in a distribution
  that is a mixture of $L$ other distributions.  Within the mixture
  setting, mixture importance sampling is a prominent example.

 If category $\ell$  has probability $\alpha_\ell$ in the motivating
  model then it would naturally get an interval of width $\alpha_\ell$.
If the RQMC method uses scrambled Sobol' points
then we can make the problem even more RQMC-friendly by instead giving
  level $\ell$ an interval of width $\beta_\ell = 2^{-\kappa_\ell}$
  for some integers $\kappa_\ell \geqslant1$ and then
  incorporating sampling ratios $\alpha_\ell/\beta_\ell$.
  These $\beta_\ell$ are negative powers of $2$ that sum to $1$.
  The number of ways to choose $\beta_1,\dots,\beta_L$
  as a function of $L$ is given by sequence A002572 in the
  online encyclopedia of integer sequences.

  Under usual assumptions on the convergence rates for RQMC
  estimates, we find that the asymptotically optimal $\beta_\ell$
  are more nearly equal than the original $\alpha_\ell$ are.
  We find some rules for making that  allocation under uncertainty
  about the appropriate RQMC rate.

  The mixing problem was studied by [1] where they use instead
  $L$ different QMC samples instead of one sample in which
  one of the variables yields $\ell$. The first QMC variable
  was used by [2] to allocate points to different processors
  in parallel computing.

\medskip
\begin{enumerate}
 \item[{[1]}] Cui, T., J. Dick, and F. Pillichshammer (2023). Quasi-Monte Carlo methods
for mixture distributions and approximated distributions via piecewise linear
interpolation. Technical report, arXiv:2304.14786 
\item[{[2]}]
Keller, A. and L. Gr¨unschlo{\ss} (2012). Parallel quasi-Monte Carlo integration by
partitioning low-discrepancy sequences. In Monte Carlo and Quasi-Monte
Carlo Methods 2010, pp. 487--498. Springer.
\end{enumerate}

\end{talk}

\begin{talk}
  {QMC confidence intervals using quantiles of randomized nets}% [1] talk title
  {Zexin Pan}% [2] speaker name
  {Johann Radon Institute for Computational and Applied Mathematics}% [3] affiliations
  {zexin.pan@oeaw.ac.at}% [4] email
  {}% [5] coauthors
  
    


The median of linearly scrambled digital net estimates has been shown to converge to the target integral value at nearly the optimal rate across various function spaces [1]. In this talk, we explore how quantiles of these estimates can be used to construct confidence intervals for the target integral. In particular, we demonstrate that as the sample size increases, the error distribution becomes increasingly symmetric, and we quantify the rate of this symmetrization for a class of smooth integrands.
            
\medskip


\begin{enumerate}
 \item[{[1]}] Pan, Zexin. (2025). Automatic optimal-rate convergence of randomized nets using median-of-means. {\it  Mathematics of Computation}, to appear.
\end{enumerate}

\end{talk}

\begin{talk}
  {Quasi-uniform quasi-Monte Carlo lattice point sets}% [1] talk title
  {Kosuke Suzuki}% [2] speaker name
  {Yamagata University}% [3] affiliations
  {kosuke-suzuki@sci.kj.yamagata-u.ac.jp}% [4] email
  {Josef Dick, Takashi Goda, Gerhard Larcher, Friedrich Pillichshammer}% [5] coauthors
  {QMC and Applications}% [6] special session. Leave this field empty for contributed talks. 
    
   

The discrepancy of a point set quantifies how well the points are distributed, with low-discrepancy point sets demonstrating exceptional uniform distribution properties. Such sets are integral to quasi-Monte Carlo methods, which approximate integrals over the unit cube for integrands of bounded variation. In contrast, quasi-uniform point sets are characterized by optimal separation and covering radii, making them well-suited for applications such as radial basis function approximation. This paper explores the quasi-uniformity properties of quasi-Monte Carlo point sets constructed from lattices. Specifically, we analyze rank-1 lattice point sets, Fibonacci lattice point sets, Frolov point sets, and $(n \boldsymbol{\alpha})$-sequences, providing insights into their potential for use in applications that require both low-discrepancy and quasi-uniform distribution. As an example, we show that the $(n \boldsymbol{\alpha})$-sequence with $\alpha_j = 2^{j/(d+1)}$ for $j \in \{1, 2, \ldots, d\}$ is quasi-uniform and has low-discrepancy.


\end{talk}
